\label{sec:legal}

\subsection{Origins of the Communities of Interest Doctrine}

The constitutional foundation for population equality in legislative districts was
established in \textit{Reynolds v.\ Sims}, 377 U.S.\ 533 (1964)~\citep{reynolds1964}.
Chief Justice Warren's majority opinion held that the Equal Protection Clause requires
that seats in both houses of a state legislature be apportioned on a population basis.
But the Court also acknowledged that rigid adherence to mathematical equality would
conflict with legitimate state policies of respecting existing political subdivisions:
\begin{quote}
  ``A State may legitimately desire to maintain the integrity of various political
  subdivisions, insofar as possible, and provide for compact districts of contiguous
  territory in apportioning its legislature.''
\end{quote}
This passage planted the seed of the communities of interest doctrine.
Political subdivisions---counties, cities, townships---are recognized as legitimate
organizing units not only because they are administratively convenient but because
they represent communities with shared interests in local governance.
The doctrine later generalized beyond formal subdivisions to any group of residents
who share economic, social, cultural, or civic interests.

\subsection{Shaw v.\ Reno and the Neutrality Requirement}

\textit{Shaw v.\ Reno}, 509 U.S.\ 630 (1993)~\citep{shaw1993}, created the most
important constraint on community-based redistricting: race cannot be the predominant
factor in drawing district lines.
Justice O'Connor's majority opinion held that a district shaped primarily to aggregate
racial groups is constitutionally suspect and triggers strict scrutiny under the
Equal Protection Clause, even if the resulting district complies with the Voting Rights
Act.
The Court in \textit{Miller v.\ Johnson}, 515 U.S.\ 900 (1995)~\citep{miller1995},
clarified that the racial predominance test applies whenever race is the ``overriding,
predominant force'' in the line-drawing process, not merely when districts have
irregular shapes.

The community character framework is designed to satisfy the Shaw-Miller neutrality
requirement.
The character vector $\mathbf{c}(t)$ contains no racial composition data.
The signals in M.1--M.7---occupational mix, land use category, housing tenure,
commuting orientation, topographic zone, administrative zone co-membership,
and transit accessibility---are all geographic and socioeconomic in nature.
They may correlate with racial composition, as virtually any geographic characteristic
does, but the correlation is incidental to the signal's purpose, not definitional to it.

A plan drawn with community character weights can be defended on the following grounds:
the algorithm is unaware of the racial composition of any tract; the weights encode
measurable, publicly documented geographic and administrative features; and the formula
is deterministic and verifiable by any party with access to the underlying data.
No human judgment about the racial character of a neighborhood enters the computation.

\subsection{State Constitutional Criteria}

Four states with independent redistricting commissions provide the clearest statutory
expressions of the communities of interest doctrine.

\paragraph{California.}
Proposition 11 (2008), which established the Citizens Redistricting Commission,
requires that after compliance with federal law, the Commission shall adopt a
plan that ``shall favor keeping together communities of interest''
(Cal.\ Gov.\ Code \S 8252(c), as enacted by Proposition 11 (2008)).
A communities of interest is defined as ``a contiguous population which shares common
social and economic interests that would be best served by common representation.''
The definition expressly excludes party affiliation and voting history.

\paragraph{Colorado.}
Article V, Section 44 of the Colorado Constitution directs the Independent Legislative
Redistricting Commission to preserve ``communities of interest, including but not
limited to ethnic, cultural, economic, trade area, geographic, and demographic
factors'' as a mandatory criterion, ranked after equal population, VRA compliance,
and contiguity.
The enumeration of ``trade area'' and ``geographic factors'' maps directly to the
M.4 commuting shed signal and the M.5 topographic signal.

\paragraph{Arizona.}
The Arizona Independent Redistricting Commission, established by Proposition 106
(2000), must give ``communities of interest'' weight as one of six criteria
(Ariz.\ Const.\ Art.\ IV, Pt.\ 2, \S1(14)).
The Arizona Commission has interpreted the criterion to include ``geographic, social,
cultural, ethnic, and economic interests,'' consistent with the M-track taxonomy.

\paragraph{Washington.}
RCW 44.05.090 instructs the Washington State Redistricting Commission to consider
``communities of interest including racial and language minority groups.''
This formulation is notable for explicitly including racial and language minorities
alongside other community types, suggesting that community membership---rather than
racial classification as such---is the organizing concept.

\subsection{The Property Tax Fiscal Bond}

Among the legally recognized instantiations of community of interest, the property
tax fiscal bond is the strongest and most judicially familiar.
Two households in the same school district, county subdivision, or municipal utility
district do not merely share a geographic location; they are legally bound together
as co-owners of a fiscal commons.
They vote on the same bond measures, pay taxes to fund the same schools and fire
stations, and receive services from the same institutions.

The Supreme Court recognized the civic significance of school district boundaries
in \textit{Milliken v.\ Bradley}, 418 U.S.\ 717 (1974)~\citep{milliken1974},
which declined to impose cross-district school desegregation remedies absent
evidence that district lines were drawn for segregatory purposes.
The Court's analysis treated school district boundaries as presumptively legitimate
community boundaries, reflecting genuine differences in local governance, tax policy,
and civic identity.
\textit{Reynolds} itself acknowledged that ``the integrity of various political
subdivisions'' is a legitimate state interest that may justify minor deviations
from strict population equality.

The M.6 signal (administrative zone co-membership) operationalizes this doctrine
by treating shared school district, fire district, water district, electric utility
service territory, and county subdivision membership as direct measures of civic
co-membership.
Tracts that share all available zone types are, in the most literal administrative
sense, members of the same community; the zone co-membership score encodes this fact
as a continuous edge weight modifier.

\subsection{State Constitutional Authority for Communities of Interest}

The communities-of-interest criterion is not merely federal — it is codified or
judicially recognized in more than 30 state constitutions and redistricting statutes.
Three landmark examples:

\textbf{California Proposition 11 (2008) and Proposition 20 (2010).}
The California Citizens Redistricting Commission Act requires that districts
``respect the geographic integrity of any city, county, city and county, local
neighborhood, or local community of interest, to the extent possible, without
unduly favoring or disfavoring any political party.''~\citep{ca2010prop20}
This is the most explicit statutory COI requirement in any large state, and it
applies both to congressional and state legislative redistricting.

\textbf{Colorado Amendments Y and Z (2018).}
Colorado's redistricting commission statutes direct commissioners to ``preserve
communities of interest,'' defined as ``groups of people with common social and
economic interests.''~\citep{co2018amendments}
The Colorado Supreme Court in \textit{In re Interrogatories on Senate Joint
Resolution 004} (2021) upheld this criterion and held that community of interest
preservation can justify minor deviations from strict population equality.

\textbf{Thornburg v.\ Gingles §2 language.}
The majority opinion in \textit{Thornburg v.\ Gingles}, 478 U.S.\ 30 (1986)~\citep{gingles1986}
discusses the geographical compactness of minority communities as a threshold
requirement for majority-minority district creation.
The concept of compactness as applied to a community — a group of people with
shared geographic identity — is the original VRA-era incarnation of the COI doctrine.
The M-track's community character signals operationalise this concept for any
community, not only racial minority communities.

\textbf{Summary.}
The community character framework has a three-layer legal foundation:
(1) federal constitutional authority (\textit{Reynolds v.\ Sims}, \textit{Shaw v.\ Reno});
(2) VRA §2 community-compactness language (\textit{Gingles});
(3) state constitutional and statutory COI criteria (California, Colorado, and 30+
other states with similar language).
This layered foundation means that community character weights are legally defensible
at every level of the redistricting litigation hierarchy.

\subsection{The ``Quantitative Neutral'' Defense}

An important practical virtue of the community character framework is that it provides
a \emph{quantitative neutral defense} against gerrymandering claims.
When a litigant challenges a district boundary as arbitrary or politically motivated,
the map-drawer using community character weights can respond:
\begin{itemize}
  \item The boundary follows the edge of minimum community similarity, computed
    from publicly available data.
  \item The formula is $w(u,v) = w_{\mathrm{base}} \times \mathrm{sim}(\mathbf{c}_u, \mathbf{c}_v)$,
    deterministic and fully disclosed.
  \item No map-drawer chose which communities to favor; the algorithm chose the
    geometrically cheapest cut, which is the cut between dissimilar communities.
  \item The formula can be replicated by any party with access to the same public data.
\end{itemize}
This is the redistricting analogue of the quantitative defenses that courts have
found persuasive in employment discrimination cases: a facially neutral criterion,
applied consistently, documented in advance, and verifiable by all parties,
is presumptively lawful even if it produces disparate geographic effects.
